\section{Global supports and intermediate transport}
\label{sec:geometry}

We prove the global supporting property by continuing the ordinary
subgradient graph from short times to every time below one.
The nonconjugacy result of \cref{sec:conjugacy} handles a possible
conjugate first cut. A local injectivity argument and covering theory
then exclude a second minimizing geodesic at that scale.

We retain the branch Hessian $A_x(w)$ and minimizing domain $G_x$
from \cref{sec:preliminaries}, and identify tangent vectors and covectors
by the metric except when explicitly using coordinate covectors.
Put $D=\operatorname{diam}M$ and
\[
 G_x=\{p\in T_xM:d(x,\exp_xp)=|p|\}=\overline{I(x)},
 \qquad c_t=\frac ct,
 \qquad Q_tf(z)=\min_x\{f(x)+c_t(x,z)\}.
\]
The total set of minimizing vectors $\{(x,p):p\in G_x\}$ is compact.
The total open injectivity domain is open, and it contains a fixed
neighborhood of the zero section.  We use the usual cut-point alternative:
a minimizing vector on its boundary either has a conjugate endpoint or
has a second minimizing geodesic to the same endpoint.

A \emph{potential} means a function represented by a continuous datum $v$ as
\[
 u(x)=\max_y\{-c(x,y)-v(y)\}.
\]
In \cref{geo:prelim}, the active logs may be defined using any such
representing datum; duality of that datum is not required. When a dual
pair is needed, we may replace it by
$v(y)=\max_x\{-c(x,y)-u(x)\}$; then $(u,v)$ is a dual pair and
$u(x)+c(x,y)+v(y)\geq0$.  Additive constants play no role below.

A finite maximum $u(x)=\max_i\{h_i-c(x,y_i)\}$ is included in this
definition after dual completion. Indeed, if
$v(y)=\max_x\{-c(x,y)-u(x)\}$, then $v(y_i)\leq-h_i$ and
$\max_y\{-c(x,y)-v(y)\}\leq u(x)$. An index active at $x$ gives the
opposite inequality. This identifies the functions, without asserting
that dual completion leaves the set of active targets unchanged.

\subsection{Supports, divided action, and subgradients}

For a semiconvex function, $\partial u(x)$ denotes its ordinary
subdifferential, with vectors identified by the metric.  Equivalently,
$p\in\partial u(x)$ if, in normal coordinates at $x$,
\[
 u(\exp_xh)\geq u(x)+\langle p,h\rangle+o(|h|).
\]
For a fixed uniform semiconvexity constant, the remainder in this
inequality can instead be bounded below by a negative quadratic term.

\begin{lemma}[Uniform supports and active logs]
\label{geo:prelim}
There is a geometric constant $C_g$ such that $c(\cdot,y)$ and
$c(x,\cdot)$ are uniformly semiconcave, with constant $C_g$, including
at cut points.  Every potential is $D$-Lipschitz and uniformly
semiconvex.  If
\[
 \mathcal V_u(x)=\{p\in G_x:
                 u(x)+c(x,\exp_xp)+v(\exp_xp)=0\},
\]
then
\begin{equation}
 \label{geo:active-hull}
 \partial u(x)=\operatorname{conv}\mathcal V_u(x).
\end{equation}
The total active-log graph
$\{(x,p):p\in\mathcal V_u(x)\}$ is compact. Every member of the convex
hull in \eqref{geo:active-hull} has a representation with positive weights using at most $n+1$
affinely independent members of $\mathcal V_u(x)$.
\end{lemma}

\begin{proof}
The divided triangle inequality is
\begin{equation}
 \label{geo:divided-triangle}
 c_{a+b}(z,y)\leq c_a(z,w)+c_b(w,y),\qquad a,b>0.
\end{equation}
It follows by concatenating constant-speed minimizing curves and using
the energy characterization of squared distance. This shortening
inequality is also used in \cite[Claim~2.4]{CMS2001}.
Equality holds when
$w$ divides a minimizing geodesic in the time ratio $a:b$.

In particular, if $p\in G_x$, $y=\exp_xp$, and $y_s=\exp_x(sp)$,
$0<s<1$, then
\begin{equation}
 \label{geo:split-support}
 c(z,y)\leq c_s(z,y_s)+\frac{1-s}{2}|p|^2,
 \qquad c_s(x,y_s)=\frac{s}{2}|p|^2.
\end{equation}
The right side is a smooth upper support near $x$, because $sp\in I(x)$.
Taking $s=1/2$, the relevant half-vectors form a compact subset of the
open injectivity domain.  The supports therefore have a common
neighborhood size and a common upper Hessian bound.  This proves
uniform semiconcavity of the cost; symmetry handles its other variable.
Also
$|c(x,y)-c(x',y)|\leq Dd(x,x')$.  Suprema of the corresponding
negative costs consequently give the asserted Lipschitz and
semiconvexity bounds for $u$.

For completeness, the directional derivative at $x$ is
\begin{equation}
 \label{geo:directional-support}
 u'(x;h)=\max_{p\in\mathcal V_u(x)}\langle p,h\rangle.
\end{equation}
The lower inequality follows from \eqref{geo:split-support} for every
active log.  To obtain the upper inequality, choose an active endpoint
at $x_r=\exp_x(rh)$ and a minimizing log there.  Compactness gives,
along any convergent subsequence, an active endpoint and a minimizing
log at $x$.  Applying the uniform upper cost support at $x_r$ to the
point $x$ bounds
$[u(x_r)-u(x)]/r$ above by the corresponding log paired with $h$,
up to $O(r)$.  This proves \eqref{geo:directional-support}.
For a semiconvex function, the ordinary subdifferential is the set of
linear forms bounded above by every directional derivative: this
follows by subtracting the fixed negative quadratic term and applying
the supporting-hyperplane characterization of a convex function.
Thus \eqref{geo:directional-support} gives \eqref{geo:active-hull}.
The contact equality defines a closed subset of the compact total
minimizing-vector set. Thus the total active-log graph is compact,
and so is each of its fibers.  Carath\'eodory's Theorem gives the finite
representation.  Choose one with the fewest positive weights: an
affine dependence would allow the weights to be varied, preserving
their sum and barycenter, until one vanished.  Thus its active
vertices are affinely independent.
\end{proof}

For a fixed base point $x$, write
\[
 H_s(p)=\frac1s A_x(sp).
\]
This is the source Hessian of the divided cost $c_s$ on the branch
specified by $sp$. We use it only when $sp$ is nonconjugate; at a cut
point the branch is the one defined in \cref{eq:branch-hessian}.
The base point $x$ will be clear from the velocity $p$.

\begin{lemma}[Radial divided-action identity]
\label{geo:radial-index}
Let $p\in G_x$, and let $0<t<s<1$.  Along
$\gamma(r)=\exp_x(rp)$, let $S_p(r)$ be the Jacobi tensor with
$S_p(0)=0$ and $D_rS_p(0)=\Id$, using parallel orthonormal frames.
Then
\begin{equation}
 \label{geo:radial-derivative}
 \frac{d}{dr}H_r(p)=-S_p(r)^{-1}S_p(r)^{-T},
 \qquad
 H_t(p)-H_s(p)\geq\kappa_g(s-t)\Id
\end{equation}
for a constant $\kappa_g>0$ depending only on the manifold.
Furthermore,
\begin{equation}
 \label{geo:radial-identity}
 A_x(w)w=w
\end{equation}
on every smooth minimizing branch.
\end{lemma}

\begin{proof}
The Hessian $H_r(p)$ is the index form of the Jacobi field having
initial value $\xi$ and final value zero on the time interval $[0,r]$.
Let $C_p$ be the other fundamental Jacobi tensor,
$C_p(0)=\Id$, $D_rC_p(0)=0$.  The field in question is
\[
 J(a)=C_p(a)\xi-S_p(a)S_p(r)^{-1}C_p(r)\xi.
\]
Integration by parts identifies its index with
$\langle S_p(r)^{-1}C_p(r)\xi,\xi\rangle$.
The Wronskian identities imply
\[
 C_p'(r)-S_p'(r)S_p(r)^{-1}C_p(r)=-S_p(r)^{-T}.
\]
Differentiating $S_p^{-1}C_p$ gives
\eqref{geo:radial-derivative}.  Jacobi evolution on the compact set
$|p|\leq D$, $0\leq r\leq1$, bounds $\|S_p(r)\|$ above by a
geometric constant.  Hence
$S_p(r)^{-1}S_p(r)^{-T}\geq\kappa_g\Id$ before conjugacy,
which proves the second assertion by integration.  Finally, for a
smooth distance branch $r$,
$\Hess(r^2/2)=r\,\Hess r+dr\otimes dr$ and
$\Hess r(\grad r,\cdot)=0$.  Since its
radial direction at the initial point is $-p/|p|$, this gives
\eqref{geo:radial-identity}; the identity at $p=0$ follows by continuity.
\end{proof}

\begin{lemma}[Transverse slack]
\label{geo:slack}
Let $p_i\in G_x$, $m_i>0$, $\sum_i m_i=1$, and
\[
 p=\sum_i m_ip_i,\qquad E=\Span\{p_i-p\}.
\]
Suppose $t\operatorname{conv}\{p_i\}\subset I(x)$ and
$0<t<s<1$.  Then
\begin{equation}
 \label{geo:transverse-slack}
 \left[\frac{A_x(tp)}t-
       \sum_i m_i\frac{A_x(sp_i)}s\right]\bigg|_{E^\perp}
 \geq\kappa_g(s-t)\Id.
\end{equation}
The same conclusion holds at a limiting time $t$ if the radially
contracted hulls lie in $I(x)$ and the displayed Hessians have
nonconjugate branch limits.
\end{lemma}

\begin{proof}
For fixed $\xi\in E^\perp$, weak MTW says that
$q\mapsto A_x(tq)(\xi,\xi)$ is concave on every segment contained
in the hull: each segment direction lies in $E$ and is therefore
perpendicular to $\xi$.  Jensen's inequality therefore gives
$H_t(p)(\xi,\xi)\geq\sum_i m_iH_t(p_i)(\xi,\xi)$.
Subtract $\sum_i m_iH_s(p_i)$ and use
Lemma~\ref{geo:radial-index} on each original minimizing ray.
The limiting assertion follows by applying this argument to the
contracted hull and passing to the specified smooth branch limits.
In particular, it does not require $sp$ to be minimizing.
\end{proof}

\subsection{The subgradient graph and its projections}

The short-time construction is a curved version of proximal minimization
and quadratic inf-convolution; compare \cite[Sections~3--5 and~7]{Moreau1965}.
We give the local estimates explicitly, since the subsequent continuation
must use only geometric and semiconvexity bounds.
For any Lipschitz semiconvex function $f$, its ordinary subgradient
graph is compact: the subgradients are bounded by its Lipschitz constant,
and the uniform local quadratic subgradient inequalities make the graph
closed under convergence of both the base point and the vector.

\begin{lemma}[Short time for a semiconvex datum]
\label{geo:short-time}
Let $f$ be Lipschitz and semiconvex.  For sufficiently small $t>0$,
depending only on its Lipschitz and semiconvexity bounds and on the
manifold, the map
\[
 F_t:\Gamma_f\longrightarrow M,\qquad
 \Gamma_f=\{(x,p):p\in\partial f(x)\},\qquad
 F_t(x,p)=\exp_x(tp),
\]
is a homeomorphism.  Its inverse is locally Lipschitz, as a map into
the tangent bundle.  The minimizing pole of $Q_tf$ is unique,
$Q_tf\in C^{1,1}$, and
\[
 \grad Q_tf(z)=-\frac1t\log_zx_t(z).
\]
The compact graph $\Gamma_f$ is an $n$-dimensional Lipschitz manifold.
\end{lemma}

\begin{proof}
If $L$ is a Lipschitz bound, a global minimizing pole satisfies
\[
 \frac{d(x,z)^2}{2t}\leq f(z)-f(x)\leq Ld(x,z),
 \qquad d(x,z)\leq2Lt.
\]
Every graph vector has length at most $L$, so its projection lies at
distance at most $Lt$ from its pole.  For small $t$, all these points
lie in a common normal coordinate ball.  On a slightly larger ball,
$D^2_{xx}c(x,z)$ is uniformly positive definite, whereas $f$ has a
fixed lower Hessian bound.  Thus $f+c_t(\cdot,z)$ is strongly convex
there.  Its global minimizer is unique, and every graph point projecting
to $z$ is the same minimizer by its stationarity equation.  Conversely
the minimizer supplies the unique log divided by $t$ as a subgradient.
This proves bijectivity; compactness proves that the inverse is
continuous.

Here is also the required quantitative local assertion. In fixed
coordinates write $p_i$ as coordinate covectors. Semiconvexity gives
\[
 \langle p_1-p_2,x_1-x_2\rangle\geq-K|x_1-x_2|^2.
\]
Insert $p_i=-D_xc(x_i,z_i)/t$.  On the above ball,
$D^2_{xx}c\geq a\Id$ and $|D^2_{xz}c|\leq B$, uniformly.  Hence
\[
 (a-tK)|x_1-x_2|\leq B|z_1-z_2|.
\]
For small $t$ this bounds the pole map in Lipschitz norm; its momentum
is then Lipschitz by the smooth stationarity equation.  Comparing
the minimum at two nearby parameters with each other's minimizing
poles gives the displayed gradient formula.  Its right side is
Lipschitz, proving $C^{1,1}$ regularity.  Since $F_t$ itself is the
restriction of a smooth map on the tangent bundle, these charts are
bi-Lipschitz.
\end{proof}

\begin{lemma}[Capture of every minimizing velocity]
\label{geo:minimizer-capture}
Let $u$ be a potential, let $t>0$, and let $z\in M$.
If $x$ is a global minimizer of $u+c_t(\cdot,z)$, then
\[
 \frac wt\in\partial u(x)\qquad\text{for every }w\in L_x(z).
\]
Consequently, $F_t:\Gamma_u\to M$ is surjective for every $t>0$.
\end{lemma}

\begin{proof}
A minimizer exists because $u+c_t(\cdot,z)$ is continuous on compact
$M$. Fix any $w\in L_x(z)$ and any $0<\lambda<1$, and put
$z_\lambda=\exp_x(\lambda w)$. Action splitting gives the upper support
\[
 U_w(x')=\frac1{t\lambda}c(x',z_\lambda)
          +\frac1{t(1-\lambda)}c(z_\lambda,z)
 \geq c_t(x',z),
 \qquad U_w(x)=c_t(x,z).
\]
It is smooth near $x$ because $\lambda w\in I(x)$, and its gradient
there is $-w/t$. The function $u+U_w$ therefore has a local minimum
at $x$. Its first-order inequality says exactly that
$w/t\in\partial u(x)$. This argument applies separately to every
minimizing log. Choosing one of them gives a point
$(x,w/t)\in\Gamma_u$ with $F_t(x,w/t)=z$.
\end{proof}

\subsection{Local injectivity and continuation}

The short-time lemma makes $\Gamma_u$ a compact topological manifold
homeomorphic to $M$. To continue its projection, we will prove local
injectivity at a scale $t$ whenever the scaled graph has reached no
conjugate point and all earlier scaled graphs lie in the injectivity
domain. The following argument does not assume that a cut has actually
occurred at $t$.

\begin{lemma}[Local injectivity of the projection]
\label{geo:local-injectivity}
Let $u$ be a potential and $0<t<1$. Suppose that
\[
 rp\in I(x)\quad\text{for every }(x,p)\in\Gamma_u
               \text{ and }0<r<t,
\]
and that every vector $tp$, $(x,p)\in\Gamma_u$, is nonconjugate.
Then $F_t:\Gamma_u\to M$ is locally injective.
\end{lemma}

\begin{proof}
Closedness of the minimizing-vector set first gives
$tp\in G_x$ for every graph point. Suppose local injectivity fails
at $(x,p)\in\Gamma_u$. There are distinct graph points
\[
 (x_j,p_j)\longrightarrow(x,p),\qquad
 (\widetilde x_j,\widetilde p_j)\longrightarrow(x,p)
\]
such that
\[
 z_j=\exp_{x_j}(tp_j)
     =\exp_{\widetilde x_j}(t\widetilde p_j)
 \longrightarrow z=\exp_x(tp).
\]

\smallskip\noindent
\emph{A common smooth branch.}
Nonconjugacy at $(x,tp)$ supplies one smooth inverse of the exponential
pair map near $(x,z)$, with branch cost $\hat c$. Both sequences of
scaled velocities belong to this branch for large $j$. Thus $x_j$
and $\widetilde x_j$ must be distinct: equality of their base points
would force equality of their velocities in the same inverse branch.
Exchange the two points if necessary so that
\[
 \Delta_j=u(\widetilde x_j)+\frac1t\hat c(\widetilde x_j,z_j)
          -u(x_j)-\frac1t\hat c(x_j,z_j)\leq0.
\]
Write $\ell_j=\log_{x_j}\widetilde x_j$, $l_j=|\ell_j|>0$,
and $e_j=\ell_j/l_j$. These logs are unique for large $j$ since
$l_j\to0$. After passing to a subsequence in a tangent-bundle chart,
$e_j\to e\in T_xM$ with $|e|=1$. All pairs $(x',z_j)$, with $x'$
on the short geodesic joining $x_j$ to $\widetilde x_j$, lie in a
fixed compact product neighborhood on which $\hat c$ is smooth.

\smallskip\noindent
\emph{Active supports and a uniform expansion.}
Choose active representations with $n+1$ slots,
\[
 p_j=\sum_{i=1}^{n+1}m_{ji}a_{ji},\qquad
 a_{ji}\in\mathcal V_u(x_j),\qquad
 m_{ji}\geq0,\qquad \sum_i m_{ji}=1.
\]
Repetitions and zero weights are allowed. Compactness of the active
graph and of the weight simplex gives, after another subsequence,
\[
 m_{ji}\to m_i,\qquad a_{ji}\to a_i\in\mathcal V_u(x),
 \qquad p=\sum_i m_i a_i.
\]
Only these finitely many active slots enter the argument; the
representing family of targets need not be finite.

Fix $s$ with $t<s<1$. Splitting the original active rays at time $s$
gives global lower supports for $u$ touching at $x_j$, with gradient
$a_{ji}$ and Hessian $-H_s(a_{ji})$. The selected divided branch cost
has gradient $-p_j$ and Hessian $H_t(p_j)$. Taylor expansion along
$\ell_j$ therefore yields, for every slot,
\begin{equation}
 \Delta_j\geq
 l_j\langle a_{ji}-p_j,e_j\rangle
 +\frac{l_j^2}{2}
   [H_t(p_j)-H_s(a_{ji})](e_j,e_j)
 +\epsilon_{ji}l_j^2,
 \qquad \max_i|\epsilon_{ji}|\longrightarrow0.
 \label{geo:collision-expansion}
\end{equation}
The remainder is uniform even for slots whose weights tend to zero.
Indeed, the branch cost has bounded third derivatives on the fixed
compact neighborhood, while the shortened active vectors $sa_{ji}$
lie in a compact subset of the open injectivity domain. The
corresponding support costs thus have one uniform third-derivative
bound. Both Taylor remainders before division by $l_j^2$ are
$O(l_j^3)$.

\smallskip\noindent
\emph{The limiting transverse directions.}
Put
\[
 d_{ji}=\langle a_{ji}-p_j,e_j\rangle,\qquad
 B_{ji}=[H_t(p_j)-H_s(a_{ji})](e_j,e_j).
\]
Every $B_{ji}$ is bounded and converges, because the first Hessian
belongs to the fixed smooth branch and the second belongs to a
uniformly shortened minimizing ray. Divide
\eqref{geo:collision-expansion} by $l_j$ and use $\Delta_j\leq0$.
The limit gives
\[
 d_i:=\langle a_i-p,e\rangle\leq0.
\]
At every finite stage the active representation gives
\begin{equation}
 \sum_i m_{ji}d_{ji}=0.
 \label{geo:collision-cancellation}
\end{equation}
Thus $\sum_i m_i d_i=0$, and $d_i=0$ whenever $m_i>0$.
The positive-weight limiting active velocities consequently lie in
one affine hyperplane perpendicular to $e$.

Let $P=\{i:m_i>0\}$ and
$E=\Span\{a_i-p:i\in P\}$. We have $e\in E^\perp$ and
$p=\sum_{i\in P}m_i a_i$. For $0<r<t$, the hull
$r\co\{a_i:i\in P\}$ lies in $r\partial u(x)\subset I(x)$.
At time $t$ all its velocities are nonconjugate by hypothesis.
The limiting transverse-slack inequality of \cref{geo:slack} therefore gives
\begin{equation}
 \delta_*:=
 [H_t(p)-\sum_{i\in P}m_iH_s(a_i)](e,e)
 \geq\kappa_g(s-t)>0.
 \label{geo:collision-gap}
\end{equation}
The later time $s$ has been applied only to the original active
velocities $a_i\in G_x$; no assertion about $sp$ or
$s\partial u(x)$ is used.

\smallskip\noindent
\emph{Cancellation before taking limits.}
Multiply \eqref{geo:collision-expansion} by $m_{ji}$ and sum over
all slots. Equation \eqref{geo:collision-cancellation} cancels the
linear terms exactly, giving
\[
 \frac{\Delta_j}{l_j^2}
 \geq\frac12\sum_i m_{ji}B_{ji}
       +\sum_i m_{ji}\epsilon_{ji}.
\]
The final sum tends to zero. The first sum tends to $\delta_*$:
positive-weight slots give \eqref{geo:collision-gap}, and zero-weight
slots disappear because their $B_{ji}$ remain bounded. No rate of
convergence of the weights is needed. We obtain
\[
 \liminf_{j\to\infty}\frac{\Delta_j}{l_j^2}
 \geq\frac{\delta_*}{2}>0,
\]
contradicting $\Delta_j\leq0$.
\end{proof}

Local injectivity has a global consequence because the projection is
homotopic to its short-time homeomorphism. We record the topological
fact in a form that does not require an orientation.

\begin{lemma}[A covering homotopic to a homeomorphism]
\label{lem:covering-homeomorphism}
Let $X$ and $N$ be nonempty connected compact topological
$n$-manifolds without boundary. A continuous locally injective map
$f:X\to N$ that is homotopic to a homeomorphism $h:X\to N$
is itself a homeomorphism.
\end{lemma}

\begin{proof}
Invariance of domain makes $f$ a local homeomorphism
\cite[Theorem~2B.3]{Hatcher2002}. Its image is open and compact,
hence also closed; connectedness of $N$ gives surjectivity.
Each fiber is compact and discrete, and therefore finite.
For $f^{-1}(y)=\{q_1,\ldots,q_k\}$, choose disjoint neighborhoods
$U_i$ on which $f$ is a homeomorphism onto a neighborhood $V_i$
of $y$. The compact set $E=X\setminus\bigcup_i U_i$ has image
missing $y$, so
\[
 W=\bigcap_i V_i\setminus f(E)
\]
is an evenly covered neighborhood of $y$. Thus $f$ is a covering.

A homotopy from $h$ to $f$ identifies their induced maps on fundamental
groups after change of base point \cite[Lemma~1.19]{Hatcher2002}.
Since $h$ is a homeomorphism, $f_*:\pi_1(X,q)\to\pi_1(N,f(q))$
is surjective. If $f(q')=f(q)$, choose a path $\gamma$ from $q$ to
$q'$. Surjectivity of $f_*$ supplies a loop $\eta$ at $q$ whose
image under $f$ is homotopic, with endpoints fixed, to $f\circ\gamma$.
Uniqueness of path lifting and the homotopy lifting property for
coverings force $\eta$ and $\gamma$ to have the same endpoint
\cite[Proposition~1.30]{Hatcher2002}. Thus $q'=q$.
A bijective local homeomorphism is a homeomorphism.
\end{proof}

\begin{theorem}[Global supporting property]
\label{geo:support}
For every potential $u$, every $x\in M$, and every
$p\in\partial u(x)$, one has $p\in G_x$ and
\begin{equation}
 \label{geo:global-support}
 u(x')\geq u(x)+c(x,\exp_xp)-c(x',\exp_xp)
 \qquad\text{for all }x'\in M.
\end{equation}
For each $0<t<1$, the projection $F_t:\Gamma_u\to M$ is a
homeomorphism; its vectors $tp$ lie in $I(x)$, and its poles are the
unique global minimizers of $Q_tu$.
\end{theorem}

\begin{proof}
The graph $\Gamma_u$ is compact by uniform semiconvexity and
Lipschitzness. Suppose that $t_0p_0\notin I(x_0)$ for some
$(x_0,p_0)\in\Gamma_u$ and some $0<t_0<1$. The continuous map
\[
 [0,t_0]\times\Gamma_u\longrightarrow TM,\qquad
 (t,(x,p))\longmapsto(x,tp)
\]
has a compact nonempty inverse image of the complement of the open
injectivity domain. Its projection onto $[0,t_0]$ has a positive
minimum $t_*$, because graph velocities are bounded and a fixed
neighborhood of the zero section lies in the injectivity domain.
Consequently all earlier scaled graph vectors lie in $I(x)$,
all vectors at $t_*$ lie in $G_x$ by closure, and an actual failure
is attained at $t_*$.

For any $p\in\partial u(x)$, represent $p$ as a finite convex
combination of active velocities $a_i\in G_x$. The hull of the
$t_*a_i$ lies in $t_*\partial u(x)\subset G_x$. Its generators are
in $I(x)$ because $t_*<1$. Corollary~\ref{geo:first-cut} therefore
makes its entire hull nonconjugate. Thus every scaled graph vector
at time $t_*$ is nonconjugate, and
\cref{geo:local-injectivity} makes $F_{t_*}$ locally injective.

Choose $\epsilon\in(0,t_*)$ small enough for
\cref{geo:short-time}. Then $\Gamma_u$ is a connected compact
topological $n$-manifold, and
\[
 (\lambda,(x,p))\longmapsto
 \exp_x\bigl(((1-\lambda)\epsilon+\lambda t_*)p\bigr),
 \qquad 0\leq\lambda\leq1,
\]
is a homotopy from the homeomorphism $F_\epsilon$ to $F_{t_*}$.
Lemma~\ref{lem:covering-homeomorphism} shows that $F_{t_*}$ is a
homeomorphism.

Fix $(x,p)\in\Gamma_u$ and set $z=F_{t_*}(x,p)$. A global minimizer
$x'$ of $u+c_{t_*}(\cdot,z)$ exists. Every minimizing log
$w\in L_{x'}(z)$ gives a graph point $(x',w/t_*)$ by
\cref{geo:minimizer-capture}; its image under $F_{t_*}$ is $z$.
Injectivity forces $x'=x$ and $w=t_*p$. Thus $x$ is the unique
global minimizing pole and $L_x(z)=\{t_*p\}$. Together with
nonconjugacy, \cref{lem:interior-characterization} gives
$t_*p\in I(x)$ for every graph point, contradicting the attained
failure. No scaled graph vector therefore reaches cut before time one.

For any fixed $0<t<1$, all hypotheses of
\cref{geo:local-injectivity} now hold: earlier scaled vectors lie in
$I(x)$, and so does the vector at $t$. The same short-time homotopy
and covering argument make $F_t$ a homeomorphism. Applying
\cref{geo:minimizer-capture} as above proves that every graph point
gives the unique global minimizing pole of $Q_tu$.

Finally, fix $(x,p)\in\Gamma_u$ and let $t\uparrow1$.
Since $tp\in I(x)$,
$d(x,\exp_xp)=|p|$. The global minimizing inequalities
\[
 u(x')+\frac1t c(x',\exp_x(tp))
 \geq u(x)+\frac1t c(x,\exp_x(tp))
 \qquad(x'\in M)
\]
pass by continuity to \eqref{geo:global-support}.
\end{proof}

\subsection{Convex sections and regular intermediate potentials}

\begin{corollary}[Convex injectivity domains]
\label{geo:convexity}
For every $x$, both $G_x$ and $I(x)$ are convex.
\end{corollary}

\begin{proof}
For $p_0,p_1\in G_x$, put $y_i=\exp_xp_i$ and consider
\[
 w(z)=\max_{i=0,1}\{c(x,y_i)-c(z,y_i)\}.
\]
Both $p_i$ belong to $\partial w(x)$, so their segment does also.
Theorem~\ref{geo:support} puts that segment in $G_x$.
Finally $I(x)$ is the interior of $G_x$: it is open, while a radial
cut-boundary vector has arbitrarily close outward radial vectors
which are not minimizing.  The interior of a convex set is convex.
\end{proof}

\begin{corollary}[Convex lifted gap sections]
\label{geo:sections}
Let $(u,v)$ be a dual pair and define
\[
 \mathcal S_x(r)=\left\{p\in G_x:
    u(x)+\frac{|p|^2}{2}+v(\exp_xp)\leq r\right\}.
\]
Every such set is compact and convex, including its ordinary and
conjugate cut endpoints.  If it is nonempty and $r\geq0$, then
\begin{equation}
 \label{geo:section-diameter}
 \operatorname{diam}\mathcal S_x(r)^2
 \leq 8\left(r+
       \osc_{p\in\mathcal S_x(r)}v(\exp_xp)\right).
\end{equation}
\end{corollary}

\begin{proof}
Write $m_p(z)=|p|^2/2-c(z,\exp_xp)$ for $p\in G_x$.
Theorem~\ref{geo:support}, applied to
$\max\{m_{p_0},m_{p_1}\}$ at $x$, gives
\[
 m_{(1-\theta)p_0+\theta p_1}(z)
 \leq\max\{m_{p_0}(z),m_{p_1}(z)\},\qquad0\leq\theta\leq1.
\]
Since
\[
 u(x)+\frac{|p|^2}{2}+v(\exp_xp)
 =\sup_z\{u(x)-u(z)+m_p(z)\},
\]
its sublevels on $G_x$ are convex.  Continuity and compactness of
$G_x$ give closedness and compactness.  For $p_0,p_1$ in a section,
their midpoint $p_m$ is also in it.  If $q(p)$ denotes the nonnegative
gap, the squared-norm identity gives
\[
 \frac{|p_0-p_1|^2}{8}
 =\frac{q(p_0)+q(p_1)}2-q(p_m)
  +v(\exp_xp_m)-\frac{v(\exp_xp_0)+v(\exp_xp_1)}2,
\]
which proves \eqref{geo:section-diameter}.
\end{proof}

\begin{proposition}[Intermediate regularity and pole control]
\label{geo:intermediate}
Let $(u,v)$ be a dual pair and $0<t<1$.  Then
\begin{equation}
 \label{geo:intermediate-duality}
 Q_tu=-Q_{1-t}v\in C^{1,1}(M).
\end{equation}
The $C^{1,1}$ seminorm is bounded solely in terms of $M,g,t$.
If $(x_t(z),p_t(z))=F_t^{-1}(z)$ and $\Phi_s$ denotes geodesic
flow on $TM$, then
\begin{equation}
 \label{geo:pole-flow}
 (x_t(z),p_t(z))
   =\Phi_{-t}(z,\grad Q_tu(z)),
 \qquad
 x_t(z)=\exp_z(-t\grad Q_tu(z)).
\end{equation}
In particular the pole and momentum are Lipschitz, with constants
depending only on $M,g,t$.

There are also $r_t,b_t>0$, depending only on $M,g,t$, such that
every projected pair $z=\exp_x(tp)$ satisfies
\begin{equation}
 \label{geo:uniform-growth}
 u(x')+c_t(x',z)
 \geq u(x)+c_t(x,z)+b_t d(x,x')^2
 \qquad\text{if }d(x,x')<r_t.
\end{equation}
These assertions use no density hypothesis.  On every compact
subinterval of $(0,1)$, their constants can be chosen uniformly.
\end{proposition}

\begin{proof}
Let $(x,p)=F_t^{-1}(z)$, set $\gamma(s)=\exp_x(sp)$, and put
$y=\gamma(1)$, $w=\dot\gamma(t)$.  Theorem~\ref{geo:support}
makes $y$ a supporting endpoint for $u$ at $x$. Since $v$ is the dual
transform, this gives
\[
 u(x)+c(x,y)+v(y)=0.
\]
Formula~\eqref{geo:divided-triangle} gives, for all $x',z',y'$,
\[
 u(x')+c_t(x',z')
 \geq -v(y')-c_{1-t}(z',y').
\]
Taking the infimum in $x'$ and the supremum in $y'$ gives
$Q_tu\geq-Q_{1-t}v$.  At $z'=z$, equality holds with $x'=x$ and
$y'=y$, because the geodesic splits in the ratio $t:(1-t)$.
This proves \eqref{geo:intermediate-duality}.

Equivalently, the function has the upper and lower supports
\begin{equation}
 \label{geo:two-sided-intermediate}
 -v(y)-c_{1-t}(z',y)
 \leq Q_tu(z')
 \leq u(x)+c_t(x,z'),
\end{equation}
both with equality at $z$.  The first is smooth there even when the
full pair $(x,y)$ is conjugate: reversing the minimizing geodesic
makes the segment from $y$ to $z$ a proper initial subsegment, and
cut symmetry gives smoothness also from $z$ to $y$.
The precise reversal identities are
\[
 \log_zx=-tw,\qquad \log_zy=(1-t)w.
\]
Thus both supports in \eqref{geo:two-sided-intermediate} have
gradient $w$.

The infimum defining $Q_tu$ preserves the uniform semiconcavity
bound of $c_t$; its representation as the supremum of the lower
supports preserves the uniform semiconvexity bound of $-c_{1-t}$.
In local coordinates it is therefore both semiconvex and semiconcave,
and hence $C^{1,1}$.  More explicitly, its almost-everywhere
Riemannian Hessian obeys
\[
 -\frac{C_g}{1-t}g\leq\Hess Q_tu\leq\frac{C_g}{t}g.
\]
The support gradients give $\grad Q_tu(z)=w$, with $|w|\leq D$.
Reversing geodesic flow proves \eqref{geo:pole-flow}.  Smooth geodesic
flow has bounded derivatives on the compact set of velocities of
length at most $D$ and times in $[-1,1]$.  The Lipschitz conclusion
therefore follows without inverting an exponential differential.

Finally choose $s=(1+t)/2$.  The single global support supplied by
$p$ can be split at $y_s=\exp_x(sp)$ to give
\[
 u(x')\geq u(x)+c_s(x,y_s)-c_s(x',y_s).
\]
After adding $c_t(x',z)-c_t(x,z)$, the right side has zero gradient
and Hessian $H_t(p)-H_s(p)\geq\kappa_g(s-t)\Id$ at $x$.
The scaled minimizing-vector families at times $t$ and $s$ are
compact subsets of the open injectivity domain.  Their smooth action
branches have uniform Taylor bounds on a common neighborhood.
Shrinking that neighborhood gives \eqref{geo:uniform-growth},
uniformly in the potential and in its graph point.
\end{proof}

\begin{remark}
\label{geo:time-one}
At time one, arbitrary potentials can have many poles with the same
endpoint; $u(x)=-c(x,y_0)$ is an example.  All uses of an inverse
projection or of $C^{1,1}$ regularization in this paper concern $0<t<1$.
Only the minimizing-vector and supporting inequalities are passed
to time one.
\end{remark}
